\section{Unbounded insertions and the optical source}
\label{sec:v80-singular-optical}
The cap in \eqref{eq:v80-incidence-local} permits summation all the
way to zero incidence. In particular the thin component is no longer
estimated by the pointwise height of another bad source.

For $0<a<5/4$ and $0<s<s_0$ put
$h_{a,s}(x)=c(x)^{-a}\one_{\{c(x)<s\}}$. The value on the
collision-null grazing set is immaterial.

\begin{theorem}[A real-frequency unbounded insertion estimate]
\label{thm:v80-unbounded-insertion}
For every fixed auxiliary band $B_*$, uniformly in the mark and the
original exact collision labels,
\begin{align}
 &m^2\int h_{a,s}(T_R^jx)
           \one_{\{K_m^{\rm c}=k,A_m=l,L_m\in J\}}\,d\nu(x)
                                                        \notag\\
 &\quad\le C_a s^{5/4-a}(|J|+B_*^{-1})
  +C_{B_*,a}(1+B_*|J|)m^2
       \min\{s^{2-a},\rho_{B_*}^{m(2-a)/4}\}.          \label{eq:v80-unbounded-local}
\end{align}
No bounded strong-space multiplication operator $M_{c^{-a}}$ is
asserted. The formula concerns the actual physical insertion; all
thresholds are summed before either the collision or band limit.
\end{theorem}
\begin{proof}
The positive layer-cake identity is
\[
 c^{-a}\one_{\{c<s\}}
  =s^{-a}\one_{\{c<s\}}
      +a\int_0^s v^{-a-1}\one_{\{c<v\}}\,dv.
\]
First truncate the integral at a positive lower endpoint. Apply
Theorem~\ref{thm:v80-superlinear-incidence} and then monotone
convergence. The leading integral is finite since $a<5/4$.
For $D\in(0,1]$ its remainder satisfies
\begin{equation}\label{eq:v80-cap-integral}
 s^{-a}\min(D,s^2)
   +a\int_0^s v^{-a-1}\min(D,v^2)\,dv
       \le C_a\min\{s^{2-a},D^{1-a/2}\}.
\end{equation}
For $s\le\sqrt D$ integrate $v^{1-a}$ directly; for larger $s$
split at $\sqrt D$. This proves the inequality for all $a<2$ and,
with $D=\rho_{B_*}^{m/2}$, proves
\eqref{eq:v80-unbounded-local}. Physical finiteness is consistent
with $\nu\{c<v\}\le v^2$ and invariance at the marked time.
No infinite sum of an uncapped spectral error is taken.
\end{proof}

Retain $d_i=\min(i,m-i)$ and define the additive optical insertion
on the original $m$-collision orbit by
\begin{equation}\label{eq:v80-depth-insertion}
 I_m(x)=1+\sum_{i=1}^{m-1}\frac{q_0^{d_i}}{c_i(x)},\qquad
 u_{n,k,m,R}(t)\,dt
       =L_\#(I_m\one_{\Pi_{n,k,m,R}}\nu_R^*).
\end{equation}
There is no product over all inverse incidences in this definition.
The source keeps its initial and terminal tests. For each finite
count it is finite and absolutely continuous: on compact regular
charts this is ordinary coarea, and the removed critical points
have zero source measure. Alternatively it is dominated by a
constant multiple of the complete inverse-incidence source in
Supplement Theorem~\ref{supp-thm:v61-inverse-incidence-density}.

\begin{corollary}[Complete depth-weighted source concentration]
\label{cor:v80-depth-source}
For every previously fixed $r>0$,
\begin{equation}\label{eq:v80-depth-source}
 \limsup_m\sup_{R,n,k}m^2\|A_r u_{n,k,m,R}\|_\infty\le C.
\end{equation}
If only the summands with $c_i<s$ are retained in $I_m$, the right
side is at most $Cs^{1/4}$. The leading constants do not depend on
$r$, $s$ or the collision count.
\end{corollary}
\begin{proof}
Apply \eqref{eq:v80-unbounded-local} with $a=1$ and sum with
weights $q_0^{d_i}$. Their sum is at most $2q_0/(1-q_0)$.
The truncated terms $c_i^{-1}\one_{\{c_i\ge s_0\}}$ and the
constant one use the inherited unweighted concentration estimate.
For each positive marked integral, dropping the endpoint tests
gives an upper bound with initial measure $\nu/c_*$ and the
\emph{same} exact coefficient $(k,n)$; its inverse-incidence
weight is retained. This upper bound does not alter the source
definition. No neighboring occupation labels
are pooled. Take $J=[t-r,t+r]$, divide by $2r$, take the collision
limsup with $B_*$ fixed, and only then let $B_*\to\infty$.
The capped remainders decay exponentially after the complete sum.
The thin version follows from the same formula at threshold $s$.
\end{proof}

\subsection{The physical normalization}
The convention of \eqref{eq:v79-optical-denominator} can also be
checked directly from the endpoint first variation. In endpoint
\emph{angles}, $F_{\alpha_0}=-Rp_0$, and hence
\begin{equation}\label{eq:v80-normalization-audit}
 \omega=|\partial_{\alpha_m}p_0|=\frac{|K_{0m}|}{R},\qquad
 d\bar\nu=\frac{\omega}{4\pi c_*}\,d\alpha_0\,d\alpha_m,
 \qquad \omega=\frac{Rc_0c_m}{y_m}.
\end{equation}
In internal \emph{arclengths} the Hessian is $D_cAD_c$.
The corner Schur complement cancels its internal incidence squares;
the two endpoint angle changes supply $R^2$, and the first-variation
conversion divides by $R$. Thus only the one factor $R$ in
\eqref{eq:v80-normalization-audit} remains. For one flight
$y_1=\ell_0$; for two flights
$y_2=\ell_0+\ell_1+2\ell_0\ell_1/(Rc_1)$. These agree with
the direct one- and two-flight Schur complements. Tangent orientation
changes conjugate the internal matrix by diagonal signs and leave
these absolute identities unchanged.

The cubic cancellation in Lemma~\ref{lem:v79-stationary-cubic}
uses these same mixed coordinates. For an internal vector $\zeta$,
differentiated stationarity gives
$D^2S[(u,Dz\,u),(0,\zeta)]=0$. Every derivative of a lifted contact
direction in $D^2F$ is purely internal, so these terms vanish before
any estimate is taken. At an endpoint the lift has arclength component
$R u_i$; the circle's second and third arclength derivatives have
norms $R^{-1}$ and $R^{-2}$. This accounts for the factors in
$C_3$ and $T_*$ without treating the endpoint angles as arclengths.

\subsection{All incidences in the noncritical optical estimate}
Let $g=\nabla_qF$ and let $\tau(x)$ be the forward lifetime of the
canonical field. For $a>0$, a superscript $a$ means restriction of
an \emph{initial} source to $|g(x)|\ge a$. Define
$O_h^{\varepsilon,a}$ by this restriction in
\eqref{eq:v79-optical-loss}. No internal incidence corridor is imposed.

A finite true exit is called grazing here when at least one internal
incidence tends to zero as the lifetime is approached. All contacts
have limits by the all-contact motion estimate. Assign a simultaneous
failure to this grazing class whenever an incidence vanishes; every
other finite exit belongs to the disjoint nongrazing class.
Let $E_h^{\varepsilon,\mathrm{gr},a}$ be its weighted starting mass
with $\tau\le h$ and $|g|\ge a$.

\begin{theorem}[Complete optical loss and genuine grazing exits]
\label{thm:v80-noncritical-optical}
For each $a,r>0$,
\begin{align}
 O_h^{\varepsilon,a}(t)+E_h^{\varepsilon,\mathrm{gr},a}(t)
          &\le \frac C a\int_0^h u(t+s)\,ds,
                                      \label{eq:v80-optical-before-limit}\\
 \limsup_m\sup_{R,n,k}m^2
  \|\bar O_r^{\varepsilon,a}
                +\bar E_r^{\varepsilon,\mathrm{gr},a}\|_\infty
          &\le \frac{Cr}{a}.          \label{eq:v80-optical-height}
\end{align}
The first assertion holds jointly for almost every $(t,h)$.
Constants are independent of $\varepsilon$ and of all positive
internal incidences and clearances. The estimate includes every
actual word of the original exact event with the stated initial
noncritical restriction.
\end{theorem}
\begin{proof}
The field and \eqref{eq:v79-optical-gradient} give, along a survivor,
\begin{equation}\label{eq:v80-optical-drift}
 |X\omega|\le C\frac{\omega I_m}{|g|}.
\end{equation}
The radial speed implies $|g(\Phi_sx)|\ge|g(x)|\ge a$.
For a surviving orbit the positive decrease of $\omega$ up to $h$
is bounded by the integral of $[-X\omega]_+$ over that segment.
At a grazing exit, \eqref{eq:v79-optical-denominator} and the
selected-incidence determinant bound imply $\omega(\Phi_sx)\to0$.
This last assertion needs only a finite-count bound
$\omega\le C_m c_i$; it does not sum such constants over words.
Thus the entire starting optical weight at such an exit is bounded
by the same negative-variation integral up to $\tau$. The two
source cases are disjoint. Other exits contribute zero to the left.

In a local angular chart the initial fiber density is
$\omega(x)/(4\pi c_*\chi(x))\,d\vartheta$. At time $s$, the ratio
$\omega(\Phi_sx)/\omega(x)$ is at most
$(\Lambda_\circ/\lambda_\circ)J_s(x)$. With $W^\varepsilon(x)\le1$,
coarea covariance in the \emph{target} direction bounds the integral
of the drift in \eqref{eq:v80-optical-drift} by
$Ca^{-1}u(t+s)$. The target image is a subset of the same exact
fiber. The common-kernel construction below justifies this comparison
simultaneously on all physical words and at joint almost-everywhere
levels. Tonelli proves \eqref{eq:v80-optical-before-limit}.
Averaging over $0<h<r$ gives an upper bound
$2Cr a^{-1}A_r u(t)$. Apply Corollary~\ref{cor:v80-depth-source}.
In particular the initial mass of a grazing exit has been paid by
optical drift; it has not been declared zero because its boundary
optical density vanishes.
\end{proof}

Without the initial lower bound, the same proof gives the extended
nonnegative inequality
\begin{equation}\label{eq:v80-critical-abel}
 O_h^\varepsilon(t)+E_h^{\varepsilon,\mathrm{gr}}(t)
                 \le C\int_0^h s^{-1/2}u(t+s)\,ds.
\end{equation}
It uses $|g(\Phi_sx)|^2\ge2\lambda_\circ s$.
A bound for every fixed average of $u$ is not a bound for this
singular convolution at its initial roof, uniformly in $m$.
Equation~\eqref{eq:v80-critical-abel} is therefore not used to remove
the critical restriction in \eqref{eq:v80-optical-height}.
